\markboth{APPENDICES}{APPENDICES}
\section{Appendix}
\label{sec:appendix}

\FloatBarrier
\subsection{Code Listings}\label{sec:appendix:code}

\begin{lstlisting}[caption={Server registry (simplified). Each URL is environment-overridable.}, label=lst:registry]
MCP_SERVERS = {
    "weather":  _url("weather",  8101),
    "routes":   _url("routes",   8102),
    "strava":   _url("strava",   8103),
    "garmin":   _url("garmin",   8104),
    "calendar": _url("calendar", 8105),
}
\end{lstlisting}

\begin{lstlisting}[caption={Server implementation pattern (simplified).}, label=lst:server_pattern]
mcp = FastMCP("weather", host="127.0.0.1",
              port=8101, stateless_http=True)

@mcp.tool()
def get_weather_forecast(lat: float, lon: float,
                         days: int = 7) -> dict:
    """Multi-day weather forecast for the given
    coordinates. Returns temperature, precipitation,
    wind, and UV index per day."""
    return _call_openmeteo(lat, lon, days)
\end{lstlisting}

\FloatBarrier
\subsection{Market and State-of-the-Art Comparison Tables}\label{sec:appendix:comparisons}

Table~\ref{tab:competitors} details the competitive landscape summarised in Section~\ref{sec:market:landscape}; Table~\ref{tab:sota_comparison} scores the most directly comparable prior systems against the research-gap dimensions of Section~\ref{sec:sota:summary}.

\begin{table}[ht]
\begin{center}
\begin{minipage}{\textwidth}
\caption[Competitive landscape of consumer sports analytics platforms]{Competitive landscape of consumer sports analytics platforms. No existing platform simultaneously provides multi-source integration, cross-domain reasoning, and conversational interaction.}\label{tab:competitors}
\begin{tabularx}{\textwidth}{p{2.0cm}p{2.4cm}p{2.2cm}X}
\toprule
\textbf{Platform} & \textbf{Core Capability} & \textbf{Data Sources} & \textbf{Key Limitation} \\
\midrule
Strava & Activity tracking, social, routes & Strava uploads & No recovery data, no weather, no cross-source reasoning \\
\addlinespace
Garmin Connect & Wellness, sleep, HRV, Body Battery & Garmin only & Walled garden; no Strava load, no route planning \\
\addlinespace
TrainingPeaks & Training plans, TSS/CTL/ATL & File import & Paid, no chat interface, manual import \\
\addlinespace
Intervals.icu & Load analytics, free tier & Strava, Garmin & No weather, calendar, or routes; dashboard-only \\
\addlinespace
WHOOP Coach & AI chat on biometrics & WHOOP device & Walled garden; no Strava, routes, or calendar \\
\addlinespace
Athletica.ai & Adaptive training plans & Garmin, Strava & Plan-only; no recovery chat, no route planning \\
\addlinespace
ChatGPT / Gemini & General-purpose LLM & None (no API) & No live fitness data; generic, ungrounded advice \\
\botrule
\end{tabularx}
\end{minipage}
\end{center}
\end{table}


\begin{landscape}
\begin{table}[p]
\centering
\small
\caption[Comparison of prior agentic and coaching systems]{Comparison of prior agentic and coaching systems against the research-gap dimensions identified in this section. \checkmark\ = capability demonstrated; $\times$ = not demonstrated or not present. Per-system detail is given in the text.}\label{tab:sota_comparison}
\begin{tabular}{p{3.6cm}p{3.2cm}p{2.5cm}p{3.5cm}p{1.5cm}p{2.0cm}}
\toprule
\textbf{System} & \textbf{Multi-Agent} & \textbf{Protocol} & \textbf{Live Data} & \textbf{RAG-Grounded} & \textbf{Domain} \\
\midrule
HuggingGPT\newline\citep{shen_hugginggpt_2023} & \checkmark & $\times$ & $\times$ & $\times$ & $\times$ \\
\specialrule{0.3pt}{3pt}{3pt}
AutoGen\newline\citep{wu_autogen_2023} & \checkmark & $\times$ & $\times$ & $\times$ & $\times$ \\
\specialrule{0.3pt}{3pt}{3pt}
CA-MCP\newline\citep{jayanti_enhancing_2026} & $\times$ & \checkmark & $\times$ & $\times$ & $\times$ \\
\specialrule{0.3pt}{3pt}{3pt}
GPTCoach\newline\citep{joerke_gptcoach_2024} & $\times$ & $\times$ & $\times$ & $\times$ & \checkmark \\
\specialrule{0.3pt}{3pt}{3pt}
SMART-ALEC\newline\citep{vemuri_multiagent_2021} & \checkmark & $\times$ & $\times$ & $\times$ & \checkmark \\
\specialrule{0.3pt}{3pt}{3pt}
WHOOP Coach\newline\citep{whoop_coach_2024} & $\times$ & $\times$ & $\times$ & $\times$ & \checkmark \\
\specialrule{0.3pt}{3pt}{3pt}
Athletica.ai\newline\citep{athletica_ai_2024} & $\times$ & $\times$ & \checkmark & $\times$ & \checkmark \\
\specialrule{0.3pt}{3pt}{3pt}
General-purpose LLM\newline\citep{puce_harnessing_2025} & $\times$ & $\times$ & $\times$ & $\times$ & $\times$ \\
\specialrule{0.3pt}{3pt}{3pt}
\textbf{Training Copilot}\newline\textbf{(this work)} & \checkmark & \checkmark & \checkmark & \checkmark & \checkmark \\
\botrule
\end{tabular}
\end{table}
\end{landscape}


\FloatBarrier
\subsection{MCP Server Tool Inventory}\label{sec:appendix:tools}

Tables~\ref{tab:full_tools}, \ref{tab:full_tools_b} and~\ref{tab:full_tools_c} provide the complete tool inventory across all native MCP servers deployed in Training Copilot.

\begin{table}[ht]
\begin{center}
\begin{minipage}{\textwidth}
\caption[Complete MCP tool inventory (1/3)]{Complete MCP tool inventory (part 1 of 3): weather, routes, and Strava.}\label{tab:full_tools}
\begin{tabularx}{\textwidth}{lllX}
\toprule
\textbf{Server} & \textbf{Port} & \textbf{Tool Name} & \textbf{Description} \\
\midrule
\multirow{4}{*}{Weather} & \multirow{4}{*}{8101}
  & \texttt{get\_current\_weather} & Current conditions \\
  & & \texttt{get\_weather\_forecast} & Multi-day forecast \\
  & & \texttt{get\_pollen\_levels} & Pollen concentrations \\
  & & \texttt{get\_uv\_index} & UV index with category \\
\addlinespace
\multirow{7}{*}{Routes} & \multirow{7}{*}{8102}
  & \texttt{geocode} & Place name to coordinates \\
  & & \texttt{plan\_route} & A-to-B route \\
  & & \texttt{plan\_circular\_route} & Loop from start point \\
  & & \texttt{plan\_park\_loop} & Loop inside named area \\
  & & \texttt{explore\_trails} & Trail search nearby \\
  & & \texttt{get\_elevation\_profile} & Elevation analysis \\
  & & \texttt{get\_isochrone} & Reachability polygon \\
\addlinespace
\multirow{14}{*}{Strava} & \multirow{14}{*}{8103}
  & \texttt{get\_activities} & Recent activities \\
  & & \texttt{search\_activities} & Full-history search by place/keyword \\
  & & \texttt{get\_activity\_stats} & Aggregate totals \\
  & & \texttt{get\_athlete\_profile} & Athlete profile + stats \\
  & & \texttt{get\_training\_trends} & Weekly volume trends \\
  & & \texttt{get\_personal\_bests} & Top performances \\
  & & \texttt{get\_yearly\_breakdown} & Year-over-year totals \\
  & & \texttt{get\_gear\_info} & Bike/shoe mileage \\
  & & \texttt{get\_activity\_detail} & Lap splits, HR, power \\
  & & \texttt{get\_activity\_streams} & Raw GPS streams \\
  & & \texttt{get\_training\_load} & CTL/ATL/TSB \\
  & & \texttt{delete\_activity} & Remove an activity \\
  & & \texttt{analyze\_performance\_trends} & Multi-metric trend analysis \\
  & & \texttt{compare\_activity\_to\_baseline} & Activity vs.\ personal baseline \\
\botrule
\end{tabularx}
\end{minipage}
\end{center}
\end{table}

\begin{table}[ht]
\begin{center}
\begin{minipage}{\textwidth}
\caption[Complete MCP tool inventory (2/3)]{Complete MCP tool inventory (part 2 of 3): Garmin, calendar, flythrough, and Google Maps.}\label{tab:full_tools_b}
\begin{tabularx}{\textwidth}{lllX}
\toprule
\textbf{Server} & \textbf{Port} & \textbf{Tool Name} & \textbf{Description} \\
\midrule
\multirow{13}{*}{Garmin} & \multirow{13}{*}{8104}
  & \texttt{get\_garmin\_activities} & Activity list \\
  & & \texttt{get\_garmin\_activity\_detail} & Per-lap splits \\
  & & \texttt{get\_garmin\_daily\_health} & Daily wellness summary \\
  & & \texttt{get\_garmin\_heart\_rate\_timeline} & Intraday HR timeline \\
  & & \texttt{get\_garmin\_sleep} & Sleep stages + score \\
  & & \texttt{get\_garmin\_body\_battery} & Energy metric timeline \\
  & & \texttt{get\_garmin\_hrv\_status} & HRV + readiness \\
  & & \texttt{get\_garmin\_training\_metrics} & VO\textsubscript{2}max, load, status \\
  & & \texttt{get\_garmin\_wellness\_trends} & Multi-day wellness trends \\
  & & \texttt{get\_garmin\_steps\_timeline} & Intraday step counts \\
  & & \texttt{get\_garmin\_stress\_timeline} & Intraday stress levels \\
  & & \texttt{get\_garmin\_body\_composition} & Weight, BMI, body fat \\
  & & \texttt{get\_activity\_gps\_track} & GPS track from activity \\
\addlinespace
\multirow{6}{*}{Calendar} & \multirow{6}{*}{8105}
  & \texttt{list\_calendars} & Available calendars \\
  & & \texttt{list\_events} & Events in range \\
  & & \texttt{get\_event} & Single event detail \\
  & & \texttt{create\_event} & Create event \\
  & & \texttt{update\_event} & Update event \\
  & & \texttt{delete\_event} & Delete event \\
\addlinespace
Flythrough & 8107
  & \texttt{prepare\_flythrough} & Package a 3D activity-flythrough render request \\
\addlinespace
\multirow{6}{*}{Google Maps} & \multirow{6}{*}{8108}
  & \texttt{maps\_search\_places} & Find places/POIs by text \\
  & & \texttt{maps\_search\_along\_route} & Places on a planned route \\
  & & \texttt{maps\_place\_details} & Hours, phone, website \\
  & & \texttt{maps\_geocode} & Address/name to coordinates \\
  & & \texttt{maps\_reverse\_geocode} & Coordinates to address \\
  & & \texttt{maps\_directions} & Directions and ETA (A$\to$B) \\
\botrule
\end{tabularx}
\end{minipage}
\end{center}
\end{table}

\begin{table}[ht]
\begin{center}
\begin{minipage}{\textwidth}
\caption[Complete MCP tool inventory (3/3)]{Complete MCP tool inventory (part 3 of 3): the athlete store (local store + deterministic training math; calls no upstream API).}\label{tab:full_tools_c}
\begin{tabularx}{\textwidth}{lllX}
\toprule
\textbf{Server} & \textbf{Port} & \textbf{Tool Name} & \textbf{Description} \\
\midrule
\multirow{14}{*}{Athlete} & \multirow{14}{*}{8109}
  & \texttt{get\_athlete\_overview} & Full athlete state (goal, zones, plan) \\
  & & \texttt{set\_race\_goal} & Create/update the main race goal \\
  & & \texttt{add\_milestone} & Add a B/C milestone \\
  & & \texttt{update\_milestone\_status} & Mark milestone done/missed \\
  & & \texttt{delete\_race\_goal} & Remove goal or milestone \\
  & & \texttt{add\_timeline\_event} & Injury/illness/race/note event \\
  & & \texttt{delete\_timeline\_event} & Remove a timeline event \\
  & & \texttt{set\_athlete\_profile} & Store the athlete's age \\
  & & \texttt{compute\_zones} & Deterministic HR + pace zones \\
  & & \texttt{scaffold\_plan} & Deterministic multi-week plan skeleton \\
  & & \texttt{save\_plan} & Persist plan (guardrail-validated) \\
  & & \texttt{record\_week\_actual} & Log a week's actual volume \\
  & & \texttt{rescaffold\_plan} & Re-baseline remaining weeks \\
  & & \texttt{get\_plan} & Current plan with progress \\
\botrule
\end{tabularx}
\end{minipage}
\end{center}
\end{table}


\FloatBarrier
\subsection{Evaluation Tables}\label{sec:appendix:evaltables}

Tables~\ref{tab:scorers}--\ref{tab:usability} give the reporting structures of the two evaluation studies of Section~\ref{sec:evaluation}: the usage study's scoring dimensions and aggregate results (conversation quality plus the trace-derived operational reliability metrics), and the usability study. Result cells are placeholders (\textbf{TBD}) pending the final runs.

\begin{table}[ht]
\begin{center}
\begin{minipage}{\textwidth}
\caption[Evaluation scoring dimensions]{Evaluation scoring dimensions. Four LLM judges run on GPT-5.4-nano; one deterministic scorer inspects tool-call traces directly.}\label{tab:scorers}
\begin{tabularx}{\textwidth}{p{3.5cm}lX}
\toprule
\textbf{Scorer} & \textbf{Type} & \textbf{Assessment} \\
\midrule
Conversation Completeness & LLM judge & Were the user's questions fully and satisfactorily answered? \\
\addlinespace
User Frustration & LLM judge & Did the user become frustrated, and did the agent worsen or mitigate it? \\
\addlinespace
Safety & LLM judge & Is the content safe and free of harmful advice? \\
\addlinespace
Supportive Coaching Tone & LLM judge & Does the assistant maintain a supportive, encouraging tone throughout, even when the user pushes back? \\
\addlinespace
Grounded in Real Data & Deterministic & Did the Copilot actually use its MCP tools to fetch real data, or did it answer from parametric knowledge alone? \\
\botrule
\end{tabularx}
\end{minipage}
\end{center}
\end{table}


\begin{table}[ht]
\begin{center}
\begin{minipage}{\textwidth}
\caption[Aggregate evaluation results (placeholder)]{Aggregate evaluation results across the ten persona conversations, reported as the mean scorer pass rate (fraction of conversations the scorer marks acceptable; for \texttt{grounded\_in\_real\_data} the fraction that issued at least one successful tool call). Values are placeholders (\textbf{TBD}) to be populated from the final run before submission.}\label{tab:results}
\begin{tabularx}{\textwidth}{Xllc}
\toprule
\textbf{Scoring Dimension} & \textbf{Type} & \textbf{Triathletes} & \textbf{Overall} \\
 & & \textbf{/ Cyclists} & \textbf{(n=10)} \\
\midrule
Conversation Completeness & LLM judge & TBD / TBD & TBD \\
\addlinespace
User Frustration (resilience) & LLM judge & TBD / TBD & TBD \\
\addlinespace
Safety & LLM judge & TBD / TBD & TBD \\
\addlinespace
Supportive Coaching Tone & LLM judge & TBD / TBD & TBD \\
\addlinespace
Grounded in Real Data & Deterministic & TBD / TBD & TBD \\
\midrule
Mean turns to goal & --- & TBD / TBD & TBD \\
Mean end-to-end latency (s) & --- & TBD / TBD & TBD \\
\botrule
\end{tabularx}
\end{minipage}
\end{center}
\end{table}


\begin{table}[ht]
\begin{center}
\begin{minipage}{\textwidth}
\caption[Usability study results (placeholder)]{Usability study reporting structure: three predefined tasks, the System Usability Scale, and an open-ended questionnaire; the participants' free-form conversations are scored by the automated pipeline (Table~\ref{tab:results}). Values are placeholders (\textbf{TBD}) to be populated once the study is complete.}\label{tab:usability}
\begin{tabularx}{\textwidth}{Xll}
\toprule
\textbf{Measure} & \textbf{Instrument} & \textbf{Result} \\
\midrule
Participants ($n$) & --- & TBD \\
\addlinespace
Task 1: add the Baden Marathon Karlsruhe as a race goal & Task completion, time on task & TBD \\
Task 2: generate a training plan preparing for it & Task completion, time on task & TBD \\
Task 3: plan today's workout (route, timing, plan-consistent) & Task completion, time on task & TBD \\
\addlinespace
Perceived usability & System Usability Scale (0--100) & TBD \\
Qualitative findings & Open-ended questionnaire (themes) & TBD \\
\botrule
\end{tabularx}
\end{minipage}
\end{center}
\end{table}


\FloatBarrier
\subsection{Usability Questionnaire and Aggregated Participant Statistics}\label{sec:appendix:ux}

The usability sessions (Section~\ref{sec:evaluation:usability}) are framed by a structured questionnaire administered as an online form. Its pre-task blocks capture the participant profile: \textit{demographics} (age, gender, occupation or field of study); \textit{sporting background} (weekly training frequency, sports practised, ambition, social training context, competition experience, training-plan use and source, motivation); \textit{fitness-app and wearable usage} (apps used, purposes, satisfaction, wearable ownership, data-analysis habits); and \textit{prior AI exposure} (usage frequency, self-rated familiarity, trust, usage contexts, willingness to pay). After the three tasks follow the ten-item System Usability Scale and closing questions (overall experience rating, likes, dislikes, suggested improvements, missing features). Figures~\ref{fig:ux_participants}--\ref{fig:corpus_usage} are placeholders for the aggregated results.

\begin{figure}[ht]
\centering
\fbox{\parbox{0.85\textwidth}{\centering\vspace{3cm}\textcolor{tcgrey!40}{\textit{Placeholder: aggregated participant charts -- age distribution, gender, occupation / field of study; weekly training frequency, sports practised, ambition (1--5), social training context, competition experience, and training-plan use and source.}}\vspace{3cm}}}
\caption[Participant demographics and sporting background (placeholder)]{Aggregated participant statistics from the intake questionnaire: demographics and sporting background. Charts to be generated once the usability study is complete.}\label{fig:ux_participants}
\end{figure}


\begin{figure}[ht]
\centering
\fbox{\parbox{0.85\textwidth}{\centering\vspace{3cm}\textcolor{tcgrey!40}{\textit{Placeholder: aggregated technology-profile charts -- fitness apps used (Strava, Garmin Connect, \ldots) and satisfaction (1--5), wearable ownership and data-analysis habits; AI usage frequency, self-rated familiarity (1--5), trust (1--5), and usage contexts.}}\vspace{3cm}}}
\caption[Participant technology and AI profile (placeholder)]{Aggregated participant statistics from the intake questionnaire: fitness-app and wearable usage, and prior AI exposure (frequency, familiarity, trust). Charts to be generated once the usability study is complete.}\label{fig:ux_tech}
\end{figure}


\begin{figure}[ht]
\centering
\fbox{\parbox{0.85\textwidth}{\centering\vspace{3cm}\textcolor{tcgrey!40}{\textit{Placeholder: SUS results -- per-item response distributions across the ten SUS statements, the aggregate SUS score (0--100), and the overall experience rating (1--5).}}\vspace{3cm}}}
\caption[System Usability Scale results (placeholder)]{System Usability Scale results: per-item distributions, aggregate score, and the overall experience rating from the closing questions. Charts to be generated once the usability study is complete.}\label{fig:ux_sus}
\end{figure}


\begin{figure}[ht]
\centering
\fbox{\parbox{0.85\textwidth}{\centering\vspace{3cm}\textcolor{tcgrey!40}{\textit{Placeholder: conversation-corpus usage charts -- conversations and messages per participant, turns per conversation, web vs.\ Telegram channel split, and specialist domains exercised.}}\vspace{3cm}}}
\caption[Conversation-corpus usage statistics (placeholder)]{Usage statistics of the real-user conversation corpus (Section~\ref{sec:evaluation:corpus}): conversations, messages, turns per conversation, channel split, and domain coverage. Charts to be generated at the end of the logging period.}\label{fig:corpus_usage}
\end{figure}


\FloatBarrier
\subsection{Environment Configuration}\label{sec:appendix:env}

Table~\ref{tab:env} lists the key environment variables used by Training Copilot.

\begin{table}[ht]
\begin{center}
\begin{minipage}{\textwidth}
\caption{Key environment variables.}\label{tab:env}
\begin{tabularx}{\textwidth}{lp{2cm}X}
\toprule
\textbf{Variable} & \textbf{Default} & \textbf{Purpose} \\
\midrule
\texttt{LLM\_PROVIDER} & \texttt{openai} & LLM backend (\texttt{openai}, \texttt{openai\_official}, \texttt{gemini}) \\
\texttt{AGENT\_MODEL} & --- & Model for agent layer \\
\texttt{AGENT\_LLM\_MODEL} & --- & Override model for agents \\
\texttt{OPENAI\_BASE\_URL} & --- & OpenAI (or compatible gateway) base URL \\
\texttt{ORS\_API\_KEY} & --- & OpenRouteService key \\
\texttt{GARMIN\_EMAIL} & --- & Garmin Connect email \\
\texttt{CLIENT\_ID} & --- & Strava OAuth client ID \\
\texttt{MLFLOW\_TRACKING\_URI} & \texttt{:5001} & MLflow server URL \\
\texttt{ORCHESTRATOR\_SPECIALISTS} & all & Enabled specialist agents \\
\botrule
\end{tabularx}
\end{minipage}
\end{center}
\end{table}


\FloatBarrier
\subsection{External API Overview}\label{sec:appendix:apis}

Table~\ref{tab:api_overview} lists every external service Training Copilot talks to over the network (including one-time/offline scripts, not just runtime agent tools), sorted alphabetically by service: what it provides, what access it requires, whether a vendor-native MCP server exists for it, and how we integrate it. Exact environment-variable names for each credential are listed in Table~\ref{tab:env}. Purely local, keyless components (the fitness RAG's embedding model in Section~\ref{sec:data:rag}, local Whisper transcription of Telegram voice memos in Section~\ref{sec:architecture:supporting}, and the self-hosted MLflow tracking server in Section~\ref{sec:agents:observability}) need no external account and are omitted here. The map-tile rows are a deliberate exception to the MCP-mediated pattern: they are fetched directly by the browser (or, for the Telegram bridge's static image, directly by the server), never through \texttt{ToolHost} or an agent.

\begin{landscape}
\small
\begin{longtable}{p{2.6cm}p{3.8cm}p{4.2cm}p{3.2cm}p{3.2cm}}
\caption[External API overview]{External APIs and network services used by Training Copilot (alphabetical), their access requirements, and MCP support.}\label{tab:api_overview}\\
\toprule
\textbf{Service} & \textbf{Provides} & \textbf{Access Required} & \textbf{Native MCP} & \textbf{Our Integration} \\
\midrule
\endfirsthead
\multicolumn{5}{l}{\textit{Table~\ref{tab:api_overview} continued from previous page}} \\
\toprule
\textbf{Service} & \textbf{Provides} & \textbf{Access Required} & \textbf{Native MCP} & \textbf{Our Integration} \\
\midrule
\endhead
\midrule
\multicolumn{5}{r}{\textit{continued on next page}} \\
\endfoot
\botrule
\endlastfoot
AWS Terrain Tiles (Terrarium) & Elevation DEM tiles for 3D terrain in the flythrough & None (public S3 Open Data bucket) & None & None (browser-fetched, outside MCP) \\
\addlinespace
CartoDB Dark Matter & Dark basemap style/tiles for the flythrough's non-satellite mode and the dashboard's dark 2D maps & None (public style/tile endpoint) & None & None (browser-fetched, outside MCP) \\
\addlinespace
ESRI World Imagery & Satellite basemap tiles for the 3D flythrough & None (public tile server) & None & None (browser-fetched, outside MCP) \\
\addlinespace
Garmin Connect & Sleep, HRV, Body Battery, stress, VO\textsubscript{2}max & Account credentials + MFA on first login (session cached); no key & None (no public API; reverse-engineered via an open-source client library) & Custom MCP server, port 8104 \\
\addlinespace
Google Calendar & Event read/write, free-slot lookup & OAuth\,2.0, or a per-request bearer token in multi-tenant mode & Community servers exist; none used here & Custom MCP server, port 8105 \\
\addlinespace
Google Geocoding & Place name $\to$ coordinates & API key (Google Cloud project) & None & Custom MCP server, port 8102 (bundled) \\
\addlinespace
Google Maps Platform & Place/POI search, place details, along-route stops, directions, and geocoding & API key (a billing-free demo key suffices); Places, Routes, and Geocoding APIs & Archived community npm proxy only (unused); we wrote our own & Custom MCP server, port 8108 \\
\addlinespace
LLM: Google Gemini & Chat completion; free-tier alternate backend & API key from Google AI Studio & N/A (consumed directly, not MCP-wrapped) & \texttt{core/llm.py}, native Gemini client \\
\addlinespace
LLM: KIT gateway & Chat completion; OpenAI-compatible alternative (tried, too slow, not used) & API key from a KIT-hosted gateway: the KIT AI Toolbox (\texttt{ki-toolbox.scc.kit.edu}) or the DSI Experimente AI Gateway (\texttt{ai-gateway.dsi-experimente.de}) & N/A (consumed directly, not MCP-wrapped) & \texttt{core/llm.py}, OpenAI-compatible client \\
\addlinespace
LLM: OpenAI (official) & Chat completion; the backend actually used for every agent, specialist, and the evaluation & API key from platform.openai.com & N/A (consumed directly, not MCP-wrapped) & \texttt{core/llm.py}, same client, official base URL \\
\addlinespace
Open-Meteo & Weather forecast, UV index, and air-quality/pollen data (two sibling endpoints) & None & None & Custom MCP server, port 8101 \\
\addlinespace
OpenRoute\-Service & A--B/circular routing, isochrones, elevation profiles & API key; free tier, 2{,}000 requests/day & None & Custom MCP server, port 8102 (bundled) \\
\addlinespace
OpenStreetMap (standard tiles) & Base map tiles for the interactive Chat/Routes maps and the Telegram bridge's static route image & None (public tile server; usage-policy limits apply) & None & None (browser-fetched, or server-side via \texttt{staticmap}; outside MCP) \\
\addlinespace
OSM Overpass & Boundary polygons for named areas (park loops) & None (public instance) & None & Custom MCP server, port 8102 (bundled) \\
\addlinespace
Project Gutenberg (Gutendex) & Public-domain book metadata and full text for the fitness RAG corpus (Section~\ref{sec:data:rag}) & None (public catalog API) & None & Offline build script only; not called at agent runtime \\
\addlinespace
Strava & Activities, load metrics (CTL/ATL/TSB), splits, GPS & OAuth\,2.0 (developer app + user consent); Standard Tier now requires a paid subscription~\citep{strava_api_update_2026} & Announced for 2026~\citep{strava_api_update_2026}, not yet used here & Custom MCP server, port 8103 \\
\addlinespace
Telegram (MTProto) & Messaging tool access (116 tools: messages, chats, contacts, media) & App API ID/hash plus a one-time session-string login & Yes, external project vendored unmodified (Section~\ref{sec:data:telegram}) & Bridged, port 8106 \\
\end{longtable}
\end{landscape}


\FloatBarrier
\subsection{UI Screenshots}\label{sec:appendix:screenshots}

Figure~\ref{fig:ui_grid} in Section~\ref{sec:architecture:llm} gives a thumbnail overview of all five Training Copilot tabs, plus the 3D flythrough modal and the Telegram delivery channel. Figures~\ref{fig:screenshot_chat}--\ref{fig:screenshot_settings} below show the same views at full size, each with a detailed description of what it renders and which part of the architecture it draws on.

\begin{figure}[ht]
\centering
\fbox{\parbox{0.85\textwidth}{\centering\vspace{4cm}\textcolor{tcgrey!40}{\textit{Screenshot: Chat interface for a multi-domain query (``Should I train tomorrow?''), showing the synthesised response, inline route map, and collapsible agent trace panel.}}\vspace{4cm}}}
\caption[Chat interface]{The chat interface, where a query triggers parallel specialist delegation and the response integrates recovery, weather, and calendar data with an interactive route map rendered from the full trace (Section~\ref{sec:agents:recording}).}\label{fig:screenshot_chat}
\end{figure}


\begin{figure}[ht]
\centering
\fbox{\parbox{0.85\textwidth}{\centering\vspace{4cm}\textcolor{tcgrey!40}{\textit{Screenshot: Health tab showing sleep stages, HRV trend against personal baseline, Body Battery timeline, and stress levels.}}\vspace{4cm}}}
\caption[Health tab]{The Health tab, rendering Garmin wellness data (sleep stages, HRV relative to baseline, Body Battery, and stress) directly from the Garmin MCP server (Section~\ref{sec:data:garmin}).}\label{fig:screenshot_health}
\end{figure}


\begin{figure}[ht]
\centering
\fbox{\parbox{0.85\textwidth}{\centering\vspace{4cm}\textcolor{tcgrey!40}{\textit{Screenshot: Dashboard tab showing an activity map, recent-activity list, training-load trend charts (CTL/ATL/TSB), and official Strava statistics (Year to Date / Last 4 Weeks / All Time).}}\vspace{4cm}}}
\caption[Dashboard tab]{The main Dashboard tab, combining an activity map, recent-activity list, training-load trend charts, and official Strava statistics for the selected sport filter.}\label{fig:screenshot_activity}
\end{figure}


\begin{figure}[ht]
\centering
\fbox{\parbox{0.85\textwidth}{\centering\vspace{4cm}\textcolor{tcgrey!40}{\textit{Screenshot: Coach tab showing the main race goal with its milestones, the computed HR/pace zones, and the multi-week training plan with a coming-weeks overview, long-run line, and plan-vs-actual view.}}\vspace{4cm}}}
\caption[Coach tab]{The Coach tab, the home of the structured race goal: the main (A) goal with its B/C milestones, the deterministically computed heart-rate and pace zones, and the goal-driven multi-week training plan (coming-weeks overview with per-week purpose and long-run line), all backed by the athlete server (Section~\ref{sec:data:athlete}); per-workout shortcuts plan a route for a session or create a calendar entry for it, reusing the existing route and calendar layers.}\label{fig:screenshot_coach}
\end{figure}


\begin{figure}[ht]
\centering
\fbox{\parbox{0.85\textwidth}{\centering\vspace{4cm}\textcolor{tcgrey!40}{\textit{Screenshot: cinematic camera following a GPS track over a satellite basemap, with elevation profile, speed readout, and progress bar overlaid.}}\vspace{4cm}}}
\caption[3D activity flythrough]{The 3D activity flythrough, where a MapLibre GL camera follows the GPS track with EMA-smoothed bearing and pitch over satellite or dark-themed basemaps, exportable as H.264 video via an in-browser WebCodecs path.}\label{fig:screenshot_flythrough}
\end{figure}


\begin{figure}[ht]
\centering
\fbox{\parbox{0.85\textwidth}{\centering\vspace{4cm}\textcolor{tcgrey!40}{\textit{Screenshot: Telegram conversation with the coach, showing a proactive daily check-in, an answer with a static route image, a tappable Google Maps directions link, and an attached GPX file, plus a transcribed voice-memo turn.}}\vspace{4cm}}}
\caption[Telegram coach conversation]{The Telegram delivery channel: the bridge forwards each message (typed or voice memo, transcribed locally) to the same orchestrator as the web chat, sends the answer back with the route as a static image, a tappable Google Maps directions link, and a GPX file, and delivers the scheduler's proactive check-ins (Section~\ref{sec:architecture:supporting}); every turn is mirrored into the pinned in-app Coach chat.}\label{fig:screenshot_telegram}
\end{figure}


\begin{figure}[ht]
\centering
\fbox{\parbox{0.85\textwidth}{\centering\vspace{4cm}\textcolor{tcgrey!40}{\textit{Screenshot: Sync tab showing a fetched range of Garmin activities, duplicate-detection status against Strava, and an export button for the missing ones.}}\vspace{4cm}}}
\caption[Sync tab]{The Sync tab, a two-stage Garmin$\to$Strava export tool: it fetches Garmin activities for a chosen date range, flags which are already present on Strava, and exports the remainder on request.}\label{fig:screenshot_sync}
\end{figure}


\begin{figure}[ht]
\centering
\fbox{\parbox{0.85\textwidth}{\centering\vspace{4cm}\textcolor{tcgrey!40}{\textit{Screenshot: Settings tab showing Strava OAuth connection, Garmin MFA login, Google Calendar authorisation, Telegram bridge setup, LLM-provider selection, and live MCP server status.}}\vspace{4cm}}}
\caption[Settings tab]{The Settings tab, where users connect Strava, Garmin, and Google Calendar through guided OAuth/MFA flows, configure the Telegram bridge and LLM provider, and see live MCP server connection status.}\label{fig:screenshot_settings}
\end{figure}

